\documentclass[a4paper,11pt]{article}
\usepackage[utf8]{inputenc}
\usepackage[T1]{fontenc}
\usepackage[french]{babel}
\usepackage[left=1cm,right=1cm,top=.5cm,bottom=0cm]{geometry}
\usepackage{enumitem}
\usepackage{hyperref}
\usepackage{xcolor}
\usepackage{titlesec}
\usepackage{fontawesome5}
\usepackage{lmodern}
\usepackage[normalem]{ulem}

% --- Couleurs (Bleu institutionnel) ---
\definecolor{primary}{RGB}{0, 51, 102}
\definecolor{secondary}{RGB}{80, 80, 80}

% --- Config Liens ---
\hypersetup{colorlinks=true, linkcolor=primary, filecolor=primary, urlcolor=primary}

% --- Titres ---
\titleformat{\section}{\large\bfseries\color{primary}\uppercase}{}{0em}{}[\titlerule]
\titlespacing{\section}{0pt}{8pt}{5pt}

\begin{document}

% --- EN-TÊTE ---
\begin{center}
    {\Large \textbf{Loïc Aron MBASSI EWOLO}} \\ \vspace{4pt}
    \textbf{\large Alternance Développeur Full Stack -- Vue.js / Node.js / MongoDB} \\
    \textit{\small Disponible dès septembre 2026 pour 1 an} \\ \vspace{4pt}
    \small
    \faMapMarker\ Cergy, France (Mobile IDF) \quad
    \faPhone\ (+33) 7 59 52 33 98 \quad
    \faEnvelope\ \href{mailto:loicaron.mbassiewolo@ensea.fr}{loicaron.mbassiewolo@ensea.fr} \\ \vspace{2pt}
    \faLinkedin\ \href{https://www.linkedin.com/in/loic-mbassi/}{linkedin.com/in/loic-mbassi} \quad
    \faGithub\ \href{https://github.com/Nameless0l}{github.com/Nameless0l} \quad
    \faGlobe\ \href{https://mbassiloic.tech}{mbassiloic.tech}
\end{center}

\vspace{-6pt}

% --- PROFIL ---
\noindent\small{Élève-ingénieur double diplôme ENSEA/Polytechnique Yaoundé avec \textbf{+3 ans d'expérience en développement web}. Maîtrise de JavaScript, React.js, et des bases de données NoSQL (MongoDB). Expérience en méthodologies Agile/Scrum et outils DevOps (Git, Docker, CI/CD). Autonome, curieux, et habitué au travail en équipe pluridisciplinaire.}

% --- FORMATION ---
\section{Formation}
\begin{itemize}[leftmargin=*, label=\tiny$\bullet$, nosep]
    \item \textbf{ENSEA -- École Nationale Supérieure de l'Électronique et de ses Applications} \hfill \textit{2025 -- 2027} \\
    \small{Double diplôme d'Ingénieur -- Systèmes Embarqués, Signal, IA. Cursus incluant architecture logicielle et DevOps.}
    \item \textbf{École Polytechnique de Yaoundé (1\textsuperscript{ère} école d'ingénieur CEMAC)} \hfill \textit{2021 -- 2025} \\
    \small{Diplôme d'Ingénieur de Conception (Master 2) : Génie Logiciel, POO, Design Patterns, Méthodes Agile, Bases de données.}
\end{itemize}

% --- COMPÉTENCES ---
\section{Compétences Techniques}
\begin{itemize}[leftmargin=*, nosep]
    \item \small{\textbf{JavaScript :} ES6+, TypeScript (maîtrise), Async/Await, Promises, DOM manipulation.}
    \item \small{\textbf{Frontend :} React.js (maîtrise), Vue.js (notions), Redux, HTML5, CSS3/SASS, Responsive Design.}
    \item \small{\textbf{Backend :} Node.js (Express, NestJS), APIs RESTful, PHP/Laravel, Python (Flask, FastAPI).}
    \item \small{\textbf{Bases de données :} MongoDB (pratique), MySQL, ScyllaDB (NoSQL). Modélisation et requêtes.}
    \item \small{\textbf{DevOps \& Outils :} Git/GitHub/GitLab (maîtrise), Docker, CI/CD, Jira, Tests unitaires, Documentation.}
    \item \small{\textbf{Méthodologies :} Agile/Scrum (stand-up, sprints, rétrospectives), Clean Code, Code review.}
    \item \small{\textbf{Langues :} Français (natif), Anglais (professionnel/technique -- collaboration internationale).}
\end{itemize}

% --- EXPÉRIENCES PROFESSIONNELLES ---
\section{Expériences Professionnelles}
\begin{itemize}[leftmargin=*, label={-}]

\item \textbf{Assistant de Recherche @ MILA (Institut Québécois d'IA)} \hfill \textit{Juin 2025 -- Sept 2025} \\
\textit{Remote -- Montréal, Canada} \hfill \textit{Python, PyTorch, Git, Documentation}
\vspace{-2pt}
    \begin{itemize}[leftmargin=0.4cm, nosep, label=\tiny$\bullet$]
        \item \small{Participation aux rituels d'équipe de recherche (stand-ups, présentations, revues de code).}
        \item \small{Développement de scripts Python pour la reproduction d'expériences état de l'art.}
        \item \small{Rédaction de documentation technique et rapports de recherche.}
        \item \small{Travail autonome en remote avec communication claire et régulière.}
    \end{itemize}
\vspace{3pt}

\item \textbf{Développeur Backend @ CSC (Cameroon Software Company)} \hfill \textit{Oct 2024 -- Jan 2025} \\
\textit{Yaoundé -- Secteur Fintech / Bancaire} \hfill \textit{Laravel, MySQL, Docker, REST APIs, Git}
\vspace{-2pt}
    \begin{itemize}[leftmargin=0.4cm, nosep, label=\tiny$\bullet$]
        \item \small{Conception et développement du backend d'une plateforme de transactions financières.}
        \item \small{Implémentation d'APIs RESTful documentées pour les intégrations tierces.}
        \item \small{Travail en équipe avec méthodologie Agile : sprints, démonstrations, rétrospectives.}
        \item \small{Utilisation de Git pour le versioning et la collaboration en équipe.}
    \end{itemize}
\vspace{3pt}

\item \textbf{Développeur Web Full Stack @ SOSDOCTA} \hfill \textit{Fév 2022 -- Mars 2024} \\
\textit{Remote (Belgique/Canada) -- Télémédecine} \hfill \textit{Laravel, PHP, MySQL, JavaScript}
\vspace{-2pt}
    \begin{itemize}[leftmargin=0.4cm, nosep, label=\tiny$\bullet$]
        \item \small{Développement complet d'une plateforme de télémédecine (backend + interfaces dynamiques JS).}
        \item \small{Maintenance évolutive sur 2 ans : nouvelles fonctionnalités, corrections de bugs.}
    \end{itemize}
\vspace{3pt}

\item \textbf{Développeur Frontend @ TOGETTECH-LLC} \hfill \textit{Fév 2023 -- Mai 2023} \\
\textit{Yaoundé -- Marketplace} \hfill \textit{React.js, Redux, JavaScript, HTML/CSS}
\vspace{-2pt}
    \begin{itemize}[leftmargin=0.4cm, nosep, label=\tiny$\bullet$]
        \item \small{Développement d'interfaces responsives avec composants React réutilisables et Redux.}
    \end{itemize}
\vspace{3pt}

\end{itemize}

% --- PROJETS ---
\section{Projets Techniques}
\begin{itemize}[leftmargin=*, label={-}]

\item \textbf{ASSO -- Marketplace Artisanat (Stack MERN)} \hfill \textit{2024} \\
\textit{Projet Freelance} \hfill \textit{MongoDB, Express.js, React.js, Node.js, TypeScript}
\vspace{-2pt}
    \begin{itemize}[leftmargin=0.4cm, nosep, label=\tiny$\bullet$]
        \item \small{Développement full stack d'une marketplace avec stack JavaScript moderne.}
        \item \small{Backend Node.js/Express avec APIs RESTful et base de données MongoDB.}
        \item \small{Frontend React/TypeScript et déploiement Docker avec CI/CD GitHub Actions.}
    \end{itemize}
\vspace{3pt}

\item \textbf{GoFind -- Architecture Microservices} \hfill \textit{2024} \\
\textit{Projet Académique} \hfill \textit{NestJS (Node.js), MongoDB, RabbitMQ, Docker}
\vspace{-2pt}
    \begin{itemize}[leftmargin=0.4cm, nosep, label=\tiny$\bullet$]
        \item \small{Architecture microservices avec API Gateway, MongoDB et communication asynchrone.}
        \item \small{Containerisation Docker et documentation OpenAPI. Travail en équipe Agile.}
    \end{itemize}
\vspace{3pt}

\item \textbf{Freengo -- Application Covoiturage} \hfill \textit{2024} \\
\textit{Projet Académique} \hfill \textit{Spring Boot, ScyllaDB (NoSQL), Next.js (React)}
\vspace{-2pt}
    \begin{itemize}[leftmargin=0.4cm, nosep, label=\tiny$\bullet$]
        \item \small{Backend Java et frontend Next.js avec base NoSQL haute performance.}
    \end{itemize}
\vspace{3pt}

\item \textbf{Laravel API Generator -- Outil Open Source} \hfill \textit{2024 -- Présent} \\
\textit{Projet Personnel} \hfill \textit{PHP, Laravel, Python, Git, Packagist}
\vspace{-2pt}
    \begin{itemize}[leftmargin=0.4cm, nosep, label=\tiny$\bullet$]
        \item \small{Générateur automatique d'architecture backend depuis UML. +30 étoiles GitHub.}
    \end{itemize}
\vspace{3pt}

\end{itemize}

% --- DISTINCTIONS ---
\section{Leadership \& Distinctions}
\begin{itemize}[leftmargin=*, label=\tiny$\bullet$, nosep]
    \item \small{\textbf{1\textsuperscript{er} Prix} IndabaX Africa 2025 (IA Santé) | \textbf{1\textsuperscript{er} Prix} CITS National Hackathon 2024 (Smart City).}
    \item \small{\textbf{Lead AI Cell} Polytechnique Yaoundé : animation de workshops React, Git, bonnes pratiques de développement.}
    \item \small{\textbf{Rédaction technique} : Articles Medium sur architecture logicielle, Git workflows et Laravel.}
    \item \small{\textbf{Certifications :} Python for Data Science (IBM), Introduction POO Java (EPFL).}
\end{itemize}

\end{document}
